\chapter{Evaluation}
\setstretch{1}

The evaluation of the HiSparse XLS Implementation was done in multiple phases. As mentioned previously, some testing for correctness was first done before codegen through the use of XLS' test proc infrastructure. After codegen was achieved, simulation of the generated verilog was done with handwritten top-level files and driven by Cocotb \cite{cocotb} coroutines with Verilator \cite{verilator} as the backend simulator. Surfer \cite{surfer}, an open-source waveform viewer, was used to visually analyze the performance of the design. The Cocotb-Verilator-Surfer evaluation flow was used from the first codegen achieved until full throughput streaming was realized. Finally, a small verilog driver was written alongside BRAM wrappers for memory requests to prepare the design for deployment onto an FPGA. For all evaluation phases, the single kernel, single cluster, two stream HiSparse architecture was used with input vector and output vector bank sizes of 2. %frequency optimization was not done on the FPGA, future work

\section{Simulation Setup}

% cocotb drivers written for the matrix loader, vector loader, vector buffer access unit and processing engine with customizable latency and backing memories

% custom python classes intaked regular formatted matrices and transformed them into iterables supporting HiSparses memory layout (of partition metadata in the front, and channels after in their sparse format)

% particularly when the monolithic proc implementations were split across blocking frontiers into send and recv procs, backpressure had to be kicked up from the recv proc up to the sending proc so that the sending proc knew when to stop, so combinational ready signals were also simulated in cocotb
% -- as an aside, due to cocotb timing model https://docs.cocotb.org/en/stable/timing_model.html, writes had to be ordered across timesteps to get a correct output

Input matrices, described in JSON files with metadata, were read by custom python classes to create iterables supporting HiSparse's memory layout and custom sparse matrix format. Similar iterable creation was done on the input vectors to support the packed data format. Cocotb drivers, i.e. python coroutines monitoring verilog wires for a particular module, were created with customizable latency for the matrix loader, vector loader, vector buffer access units, and processing engines to simulate a memory servicing their requests. The backing store for the drivers could be the aformentioned matrix and vector iterables for the matrix loader and vector loader drivers, or simple empty arrays to reflect the banks of the vector buffer access units and processing engines. 

When the previously monolithic proc implementations were split across blocking frontiers into send and receive procs, any ready backpressure from the receive proc must be routed upwards towards the send proc to ensure the send proc pauses iteration. Combinational ready connections were simulated within the Cocotb drivers to simulate this backpressure, which also required the ordering of reads and writes within the Cocotb driver code to ensure correctness with regards to Cocotb's timing model \cite{cocotb}. Figure~\ref{cocotbtimingmodel} displays the phases of a given Cocotb simulation time step in the timing model. To ensure a write observed up to date data, an arbitrary number of HDL Evaluation -> Values Changed loops were waited upon before sampling wire values, while other writes and reads executed after a lesser number of iterations of said loop. This difference in loop iterations established an order between wire reads and writes and ensured wire connections settled into correct values.

\begin{figure}[H]
\centering\includegraphics[alt={sample image},width=15cm]{phasestimingmodel.png} 
\caption[Cocotb Timing Model Phases~\CiteNolink{cocotb}.]{Phases of a time step in Cocotb's Timing Model~\cite{cocotb}.}\label{cocotbtimingmodel}
\end{figure}

% skid buffers were hand-written in verilog to ensure pipelined metadata was attached to returned values when ready was deasserted from the aformentioned backpressure

% To reduce the manual inputs necessary to generate each XLS module, a makefile was developed that replaced parameterized values with presets and generated it with saved codegen flags. I was not able to find a way to generate a proc with parameters from the command line so this was my solution

Whenever a receive proc exerted backpressure on upstream send procs, multiple in flight memory requests, the number determined by the latency of the simulated memory, must be saved to ensure values are not dropped once the receive proc is ready again. This backpressure was managed using coroutines alongside Cocotb's event triggers. Each transaction fired its own coroutine that only progressed when downstream modules were ready and ensured results returned in order.

Finally, to reduce the number of manual inputs necessary to codegen each XLS module, a Makefile was developed that generated the modules of the XLS HiSparse implementation with preset codegen flags.

\section{Simulation Results}

% Compare the cycle count of each version. Also talk about what difference between each version with the previous. Take a look at your git commits for inspiration as well.

Table~\ref{cycleres} displays the number of cycles used to calculate the SpMV result on the benchmark matrix and vector displayed in Figure~\ref{inputmatrix}. The first version, the Unoptimized Conservative, took 2058 cycles to compute the SpMV result. Through minor codegen tweaks such as generating with a target II of 1, its optimized variant achieved 897 cycles. This was the upper bound of the conservative framework, and the following versions required the architectural changes mentioned previously to achieve better performance. The Optimized Matrix Loader version split the matrix loader into send and receive procs, however, it resulted in no speed increase from the Optimized Conservative. Through analyzing the waveform generated during the cycle simulation, the unoptimized shuffle unit bottlenecked the performance of the design in both scenarios. Correspondingly, once the shuffle unit was optimized, a 70\% speed increase was observed. The shuffle unit is used twice within a cluster to manage both input vector and output vector bank conflicts. It also had the most complicated datapath out of all modules, requiring many activations for each new input value due to the multiple states necessary for communication between the shuffle module, shuffle core and arbiter. Thus, much of this speed increase can be attributed to the optimized shuffle unit.

Inbetween the two shuffle units is the vector buffer access unit, whose optimization saved 136 cycles. The vector buffer access unit only had a single blocking frontier and was a relatively simple state machine originally, explaining the modest speed increase. Finally, the optimization of the processing engine resulted in the final 155 cycles saved and represents the full streaming throughput design. Like the shuffle unit, the processing engine required multiple activations for each incoming payload due to the multiple states necessary for its acccumulation. Thus its streaming optimized variant saved a few more cycles than the vector buffer access unit.

\begin{figure}[H]
\centering\includegraphics[alt={sample image}, width=6.67cm]{input_matrix.png}
    \caption[Benchmark Matrix and Vector]{Matrix and vector used to benchmark each codegen capable HiSparse version. The four column partitions are highlighted in different colors, and each stream is its own shade of the partition's color.}\label{inputmatrix}
\end{figure}

\begin{table}[H]
\centering
\begin{tabular}{|l|r|r|}
\hline
Version & Cycle Count & Percent Speedup from Previous \\ \hline
Unoptimized Conservative & 2058 & n/a\\ \hline
Optimized Conservative & 897 & 129\%\\ \hline
Optimized Matrix Loader & 897 & 0\%\\ \hline % need to explain the lack of speed increase here
Optimized Shuffle Unit & 528 & 70\%\\ \hline % can talk about the SOD syncers dropping this from 608 to 528
Optimized Vector Buffer Access Unit & 392 & 35\%\\ \hline
Optimized Processing Engine & 237 & 65\% \\ \hline
\end{tabular}
\caption[Simulated Cycle Performance]{Simulated cycle performance of each codgen capable HiSparse XLS design. Versions are introduced in order of development and include the optimizations of the previous versions if applicable.}\label{cycleres}
\end{table}

% figure of the first partition of the matrix loader, show
\begin{figure}[H]
\centering\includegraphics[alt={sample image}, width=10cm]{optimized_matrix_loader.png}
    \caption[Matrix Loader Waveform Diagram]{Comparison of the streaming performance between the monolithic and split-proc versions of the matrix loader. At the top, each valid address is separated by multiple unused cycles due to the conservative state machine as well as the generated pipeline stages from the blocking frontier. At the bottom, multiple addresses are communicated without empty cycles inbetween, highlighting the full throughput streaming performance of the optimized matrix loader.}\label{optimizedmatrixloader}
\end{figure}

% talk about the ideal number of operations done for a dense matrix calculation and for this particular sparse calculation. The answer is 120, which i never achieve due to the overhead of syncing and latency of the pipeline. Larger matrices I would probably beat dense on. Perhaps I should run on a larger matrix then, which I did (a 32x32 99% sparse one) and got 727 for the fully opt

% talk about a 99% 32x32 sparse matrix taking 727 cycles which beats the trivial case

For a $M$ by $N$ matrix, there are $M$ multiplies and $M-1$ additions for each output vector element calculated in a dense matrix-vector multiplication, thus the equation for the number of operations is $N*(M+(M-1))=N*(2*M - 1)$. Using this formula, a simple implementation that performed a dense matrix-vector multiplication on the benchmark matrix and took one cycle per operation would take 120 cycles. Thus despite the HiSparse design achieving full throughput during streaming, the unoptimized initialization and cleanup-phases as well as general overhead of the HiSparse modules contributed enough latency for the SpMV calculation to exceed the cycle count of the simple, dense-assuming implementation. To realize the efficiency gains of the SpMV architecture, a larger 32x32 matrix that was ~99\% sparse was ran on the same single kernel, single cluster two stream architecture and achieved a runtime of 727 cycles. This is 177\% faster than the simple, dense-assuming runtime of 2016 cycles using the above formula, displaying the greater efficiency of the HiSparse design.

% talk about throughput difference between actual and optimized version of modules with Surfer screenshots

% perhaps have a screenshot of surfer for the entire pipeline as well just to show the overall difference? Talk about remaining ready deassertions in the pipeline likely caused by bank conflicts

\section{Hardware Setup}
% physical hardware setup
% also BRAM wrappers
% also coe files
% also block diagram setup and AXI GPIO stuff to read/write
A Pynq-Z2 \cite{pynqz2} served as the target board and its Zynq-7000 SoC the target chip for synthesizing the design. Vivado 2024.2 \cite{amdvivado} was used to generate the bitstream and hardware overlay files to initialize the design in the Processing System's (PS) Jupyter Notebook environment provided within the Pynq ecosystem. A wavlink AC1200 router was used to set up a LAN network to communicate between the Pynq and host machine.

Some verilog was written outside of the XLS environment to prepare the design for synthesis. A verilog testbench was created to drive the cluster in a similar manner as the Cocotb simulation and record the results and cycle latency of the SpMV calculation. Likewise, all memory requests were converted to support the Block RAM (BRAM) interface \cite{bram} using intermediate wrappers. As metadata must be pipelined alongside memory requests due to the frontier splitting optimization technique, these wrappers included skid buffers to both ensure values were not dropped whenever downstream modules exhibited backpressure, and that values were always associated with their correct metadata. Xilinx Coefficient files (COE) \cite{coefile} were used to prepopulate the BRAMs with the input matrix and vector data using HiSparse's sparse format.

To ensure the design was running properly, AXI GPIO IP blocks \cite{axigpio} were included in the block diagram and used to sample the values of the output vector and cycle count wire of the design. These wires were read on the PS side of the Zynq using python function calls.

\begin{figure}[H]
\centering\includegraphics[alt={sample image}, width=\linewidth]{single_cluster_rtl_block_diagram.png}
    \caption[Block Design]{A portion of the block design containing the HiSparse RTL module and its related BRAMs. Note the processing engine's dual port nature, as pipelined accumulations require concurrent read and writes.}\label{blockdiagram}
\end{figure}

\section{Hardware Results}
% floorplanning and utilization

The target frequency during synthesis was 100 Mhz. The design originally failed to meet this timing requirement and required two modificaitons along the critical path. The first involved a result draining module that performed a 32 bit modulo to accumulate results across all kernels and clusters. Since the design is a single kernel, single cluster variant with two streams, the bit width of the modulo's inputs could be easily reduced to meet timing without any practical loss of performance. However, the second offending path involved the processing engines multiply accumulate. Figure~\ref{timingreport} displays a setup timing report of an example offending path.

\begin{figure}[h]
\centering\includegraphics[alt={sample image}, width=\linewidth]{timingreport.png}
    \caption[Timing Report]{The timing report of the processing engine critical path. }\label{timingreport}
\end{figure}

Analyzing the path's timing, it arrives at the clock edge at 14.830ns while its required time is 12.749ns, consequently overshooting the requirement by 2.081ns. A substantial amount of the timing delay is the result of high fanout clock nets. However, in the middle of the timing report lies two DSP48E1 blocks \cite{dsp48} in succession. The DSP48E1 is a primitive used within the Programming Logic (PL) fabric to perform the multiply accumulate operation. As the processing engine performs a multiply accumulate, it is reasonable that vivado inferred its use. However, the DSP48E1 accepts a 25 bit and 18 bit twos complement input respectively for its multiplication operation. The original processing engine implementation performed a 32bitx32bit signed multiplication ideally within a cycle (under full throughput streaming operation). To support a higher bitwidth on both ports, multiple DSP48E1 blocks must be connected together in a hierarchal fashion, and consequently two of these DSPs end up being connected in succession. These successive DSPs contributed 6.53ns of delay, with the second DSP contributing 2.438ns. Thus allowing the matrix and vector elements to be 32 bit led to negative slack violations, and the removal of the second DSP would make the design achieve the timing requirements.

Reducing the bitwidth of the matrix and vector elements to 18 bits resulted in the design achieving the timing constraints and producing the correct output as it could fit in a single DSP48E1. Figure~\ref{singlelinemodification} displays this simple modification. Though technically one of the operands can be 25bits and still fit within the DSP, XLS did not support asymmetrical bitwidths during multiplication, hence both operands are 18 bits. This modificaiton affects performance by hurting the dynamic range of the design, though it enables its operation at 100 Mhz on the Zynq-7000 SoC.

\begin{figure}[H]
\begin{lstlisting}[label={singlelinemodification},caption={Single line modification necessary to acheive 100Mhz. Multipy inputs wrapped with s18 casts. Accumulation with the 32bit increment is still supported as the DSP48E1 supports up to 48 bit accumulation}, captionpos=b,style=mystyle]
pub proc pe_recv {
    // ...
    let incr = ((spld.matrix_val as s18) * (spld.vector_val as s18)) as s32;
    // ...
}
\end{lstlisting}
\end{figure}

% analyze the timing report. Why is there high fanout? The mult DSP itself supports 600 Mhz, why is it failing here? why is the path so long?

% do i ask chris if there is a way to break up the signed multiply into multiple stages?
% yea i do, i sent the question over

% modifications necessary to make timing at 100 mhz
% current critical paths and their analysis (its still the 16 bit multiply i think but take a look at the other paths)

\begin{figure}[H]
\centering\includegraphics[alt={sample image}, width=\linewidth]{single_cluster_rtl_synthesized_design.png}
    \caption[Design Floorplanning]{Synthesized design floorplanning. Blocks highlighted in purple represent the logic encoding the HiSparse architecture. Blocks highlighted in blue represent the BRAMs used by the HiSparse modules. All other logic (processing system, GPIO, etc) is highlighted in yellow.}\label{synthdesign}
\end{figure}

\begin{table}[H]
\centering
\begin{tabular}{|l|r|r|r|r|r|}
\hline
Site Type & Used & Fixed & Prohibited & Available & Util\% \\ \hline
Slice LUTs & 11786 & 0 & 0 & 53200 & 22.15 \\ \hline
\quad LUT as Logic & 11574 & 0 & 0 & 53200 & 21.76 \\ \hline
\quad LUT as Memory & 212 & 0 & 0 & 17400 & 1.22 \\ \hline
\quad\quad LUT as Distributed RAM & 18 & 0 & -- & -- & -- \\ \hline
\quad\quad LUT as Shift Register & 194 & 0 & -- & -- & -- \\ \hline
Slice Registers & 22380 & 0 & 0 & 106400 & 21.03 \\ \hline
\quad Register as Flip Flop & 22380 & 0 & 0 & 106400 & 21.03 \\ \hline
\quad Register as Latch & 0 & 0 & 0 & 106400 & 0.00 \\ \hline
F7 Muxes & 3 & 0 & 0 & 26600 & 0.01 \\ \hline
F8 Muxes & 1 & 0 & 0 & 13300 & $<0.01$ \\ \hline
\end{tabular}
\caption[Design Resource Utilization]{Resource utilization summary from synthesis. LUTs for logic and registers used 42.79\% of all slices on the Zynq-7000, while only ~1\% of LUTs were used as memories since the benchmarked matrix was so small. This resource utilization seems to support another cluster, kernel or more streams, though that will limit the available LUTs to be used as BRAMS and consequently the maximal input matrix and vector size.}\label{utilization}
\end{table}

